\chapter{Provider Configuration Matrix}
\label{app:providers}

Construction, authentication, and the places the provider abstraction leaks.

\begin{note}
Chapter~\ref{ch:first-agent} refers here for measurements. \S\ref{sec:measure-template}
is a recording template rather than a finished table --- the numbers depend on
your hardware, your region, and the model versions available when you run it, so a
table filled in by someone else would be actively misleading. The methodology is
fixed; the numbers are yours.
\end{note}

% ============================================================
\section{Construction}

Every listing below produces an \api{IChatClient}. What you do with it afterwards
is identical across providers --- see \S\ref{sec:seam} in
Chapter~\ref{ch:first-agent}.

\subsection{Ollama (local)}

\begin{csharp}[caption={Ollama via the first-party connector}, label={lst:prov-ollama}]
// dotnet add package Microsoft.Extensions.AI.Ollama --prerelease
using Microsoft.Extensions.AI;

IChatClient chatClient = new OllamaChatClient(
    new Uri("http://localhost:11434"),
    modelId: "llama3.2");
\end{csharp}

\begin{csharp}[caption={Ollama via OllamaSharp}, label={lst:prov-ollamasharp}]
// dotnet add package OllamaSharp
using OllamaSharp;

IChatClient chatClient = new OllamaApiClient(
    new Uri("http://localhost:11434"),
    "qwen3");
\end{csharp}

No authentication. Default port \code{11434}. Both clients implement
\api{IChatClient}, so either works; the samples in this book use the first.

\subsection{Azure OpenAI}

\begin{csharp}[caption={Azure OpenAI with credential-based authentication}, label={lst:prov-azure}]
// dotnet add package Azure.AI.OpenAI --prerelease
// dotnet add package Azure.Identity
using Azure.AI.OpenAI;
using Azure.Identity;
using Microsoft.Extensions.AI;

IChatClient chatClient = new AzureOpenAIClient(
        new Uri(config["AzureOpenAI:Endpoint"]!),
        new DefaultAzureCredential())
    .GetChatClient(config["AzureOpenAI:Deployment"]!)
    .AsIChatClient();
\end{csharp}

Note that the argument to \api{GetChatClient} is a \emph{deployment} name, not a
model name --- they are frequently different, and the resulting error when you
confuse them is unhelpful.

\api{DefaultAzureCredential} walks a chain of credential sources. On a
development machine that usually resolves to your \code{az login} session; in
Azure it resolves to a managed identity. Prefer it over
\api{AzureKeyCredential}: there is then no secret to leak.

\subsection{OpenAI}

\begin{csharp}[caption={OpenAI}, label={lst:prov-openai}]
using OpenAI;
using Microsoft.Extensions.AI;
using System.ClientModel;

IChatClient chatClient = new OpenAIClient(new ApiKeyCredential(apiKey))
    .GetChatClient(modelId)
    .AsIChatClient();
\end{csharp}

\subsection{Microsoft Foundry}

Foundry differs from the others: the project client becomes an agent directly,
with the model named at agent construction rather than at client construction.

\begin{csharp}[caption={Foundry Agents}, label={lst:prov-foundry}]
using Azure.AI.Projects;
using Azure.Identity;
using Microsoft.Agents.AI;

AIAgent agent = new AIProjectClient(new Uri(endpoint), new DefaultAzureCredential())
    .AsAIAgent(
        model: deploymentName,
        instructions: "You are a friendly assistant.");
\end{csharp}

\begin{warning}
Foundry stores conversations server-side. A serialised \api{AgentThread} from a
Foundry agent therefore contains an \emph{identifier}, not the messages --- and
cleaning up the remote thread is your responsibility, not the framework's. See
\S\ref{sec:threads}.
\end{warning}

\subsection{Anthropic}

\begin{verifybox}
\pkg{Microsoft.Agents.AI.Anthropic} exists at \mafcore{}-preview and its
description reads ``support for Anthropic Agents,'' but the construction syntax
has not been verified by the author. Check the package's own documentation before
using it. Do not assume it mirrors Listing~\ref{lst:prov-openai}.
\end{verifybox}

\subsection{Amazon Bedrock and Google Gemini}

Both are named as supported providers in Microsoft's 1.0 announcement. Neither
has a package under the \code{MicrosoftAgentFramework} account
(Appendix~\ref{app:packages}), which means they are reached through an
\api{IChatClient} implementation from elsewhere in the
\pkg{Microsoft.Extensions.AI} ecosystem rather than through an Agent
Framework-branded package.

\begin{verifybox}
Identify the correct connector package for these two before writing against them.
The pattern will be the one in Listing~\ref{lst:prov-ollama} --- construct a
provider client that implements \api{IChatClient}, then proceed identically ---
but the package name is unverified here.
\end{verifybox}

% ============================================================
\section{Function invocation}

Some \api{IChatClient} implementations require the function-invocation middleware
to be added explicitly before tool calling works:

\begin{csharp}[caption={Adding function invocation to a chat client}, label={lst:prov-funcinvoke}]
using Microsoft.Extensions.AI;

var agent = chatClient
    .AsBuilder()
    .UseFunctionInvocation()
    .Build()
    .CreateAIAgent(
        name: "ToolUser",
        instructions: "You are a useful agent.",
        tools: [AIFunctionFactory.Create(GetWeather)]);
\end{csharp}

\begin{note}
Whether this is required or already wired depends on the client. If tools are
registered, the model is capable, and nothing is ever called, this is the first
thing to check --- before you start rewriting tool descriptions. See also
Appendix~\ref{app:troubleshooting}, \S\ref{sec:tb-no-tool-call}.
\end{note}

% ============================================================
\section{Support characteristics}

What can be stated without measurement.

\begin{center}
\small
\begin{tabularx}{\linewidth}{@{}lXX@{}}
\toprule
\textbf{Provider} & \textbf{Thread storage} & \textbf{Auth} \\
\midrule
Ollama & In-memory --- serialised thread contains the messages & None \\
Azure OpenAI & In-memory & Entra credential or key \\
OpenAI & In-memory, or server-side with Responses & API key \\
Foundry Agents & Server-side --- serialised thread contains an id & Entra credential \\
Anthropic & Unverified & API key \\
\bottomrule
\end{tabularx}
\end{center}

The thread-storage column is the one with architectural consequences. It
determines the size and content of what you persist, whether conversation data
leaves your infrastructure, and whether you have remote state to clean up.

% ============================================================
\section{Measurement template}
\label{sec:measure-template}

Fill this in yourself. The exercise is in \S\ref{sec:measure} of
Chapter~\ref{ch:first-agent}.

\subsection{Method}

\begin{enumerate}[leftmargin=1.4em]
\item Ten fixed prompts: five open conversational, five requiring a specific
      structured answer. Same prompts for every provider, stored in the repo so
      the run is reproducible.
\item Twenty runs per prompt per provider. Discard the first run per provider ---
      it measures connection setup, not inference.
\item Record wall-clock time to first streamed update and to completion, and the
      token counts the provider reports.
\item Run providers in interleaved order rather than in blocks, so that network
      conditions or thermal throttling do not load onto one provider.
\item Record the machine, the model version, the region, and the date. Without
      these the numbers are not comparable to anything, including your own later
      runs.
\end{enumerate}

\subsection{Results}

\begin{center}
\small
\begin{tabularx}{\linewidth}{@{}Xrrrrr@{}}
\toprule
\textbf{Provider / model} & \textbf{TTFT p50} & \textbf{TTFT p95} & \textbf{Total p50} & \textbf{Tok in} & \textbf{Tok out} \\
\midrule
Ollama / llama3.2 & & & & & \\
Ollama / qwen3:4b & & & & & \\
Azure OpenAI / & & & & & \\
OpenAI / & & & & & \\
Foundry / & & & & & \\
\bottomrule
\end{tabularx}
\end{center}

\emph{Machine:} \dotfill \\
\emph{Date:} \dotfill \qquad \emph{Region:} \dotfill

\subsection{Tool-selection accuracy}

Complete this after Chapter~\ref{ch:tools}. Selection rate is the fraction of
twenty runs in which the correct tool was called with correct arguments.

\begin{center}
\small
\begin{tabularx}{\linewidth}{@{}Xrrr@{}}
\toprule
\textbf{Provider / model} & \textbf{Simple args} & \textbf{Nested object} & \textbf{Collection} \\
\midrule
Ollama / llama3.2 & & & \\
Ollama / qwen3:4b & & & \\
Azure OpenAI / & & & \\
OpenAI / & & & \\
\bottomrule
\end{tabularx}
\end{center}

\begin{note}
Expect the gap between local and frontier models to be widest in the rightmost
column. That gap, rather than raw latency, is usually what decides whether a
cheap provider is viable for a given workload --- and it is the number almost
nobody measures before committing.
\end{note}

\subsection{The question the numbers answer}

Not ``which provider is fastest.'' The useful question is: \emph{for which of my
own workloads would the cheapest option be sufficient?} Write the answer here in
prose once the tables are complete, because that sentence is the only part of
this appendix you will still care about in six months.
